% Section 4 - Lambdas and callbacks
% Roberto Masocco <roberto.masocco@uniroma2.it>
% April 22, 2026

% ### Lambdas and callbacks ###
\section{Lambdas and callbacks}
\graphicspath{{figs/section4/}}

% --- Lambdas and callbacks ---
\begin{frame}{Lambdas and callbacks}{Motivation}
	\justifying
	Many systems need to invoke \textbg{user-supplied code} at a point determined by the framework, not the user — a button press, a message arrival, a timer firing... on some \textbg{event}.\\
	This pattern is called a \textbg{callback}.\\
	\bigskip
	In C, callbacks are plain \textbg{function pointers}: a fixed signature, no associated state, no closure.\\
	Passing context requires an extra \texttt{void *} argument and a cast — fragile and untyped.\\
	\bigskip
	C++ addresses this with two complementary tools:
	\begin{itemize}
		\item \textbg{lambda expressions}: anonymous functions that can \textbg{capture} local variables from their surrounding scope;
		\item \texttt{std::function}: a type-erased \textbg{wrapper} that can hold \textbg{any callable} — a lambda, a function pointer — with a given signature.
	\end{itemize}
	\begin{block}{}
		\centering
		\textbf{These two features are the backbone of event-driven programming in C++.}\\
		Every ROS 2 subscriber callback and timer callback is expressed through them.
	\end{block}
\end{frame}
\begin{frame}[fragile]{Lambdas and callbacks}{Lambda syntax}
	\justifying
	A \textbg{lambda expression} defines an anonymous function inline, at the point of use, as a value.\\
	\begin{columns}
		\column{.9\textwidth}
		\begin{lstlisting}[style=codebeamer, language=C++, caption=Lambda syntax and basic examples.]
// General form: [capture](params) -> return_type { body }
auto add = [](int a, int b) { return a + b; };
std::cout << add(3, 4) << std::endl; // 7

// Lambdas can be passed directly where a callable is expected
std::vector<int> v = {5, 1, 4, 2, 3};
std::sort(v.begin(), v.end(), [](int a, int b) {
  return a < b;
});
    \end{lstlisting}
	\end{columns}
	Lambdas are \textbg{first-class values}: they can be stored in variables, passed as arguments, and returned from functions.
\end{frame}
\begin{frame}{Lambdas and callbacks}{Capture}
	\justifying
	The \textbg{capture clause} \texttt{[...]} controls which \textbg{variables} from the \textbg{enclosing scope} are visible inside the lambda body, and how.\\
  \bigskip
  \textbg{Capture by reference} requires that the captured variable \textbg{outlives} the lambda — a common source of dangling-reference bugs.
\end{frame}
\begin{frame}[fragile]{Lambdas and callbacks}{Capture}
  \begin{columns}
		\column{.9\textwidth}
		\begin{lstlisting}[style=codebeamer, language=C++, caption={Capture by value, by reference, and \texttt{this}}.]
int threshold = 10;

// Capture by value: lambda gets its own copy
auto by_val = [threshold](int x) { return x > threshold; };
// Capture by reference: lambda sees the original variable
auto by_ref = [& threshold](int x) { return x > threshold; };
threshold = 20; // by_ref will now use 20

// Inside a member function, capture [this] to access members
class Foo {
  int lim_ = 5;
  void bar() {auto baz = [this](int v) { return v < lim_; };}
};
    \end{lstlisting}
	\end{columns}
\end{frame}
\begin{frame}{Lambdas and callbacks}{\texttt{std::function} and \texttt{std::bind}}
	\justifying
	The type of a lambda is \textbg{anonymous and compiler-generated} — it cannot be written explicitly.\\
	\texttt{std::function<R(Args...)>} is the solution: a \textbg{type-erased} wrapper that can hold \textbg{any callable} with a matching signature.\\
  \bigskip
  \texttt{std::bind} instead creates a callable by \textbg{binding} some (or all) arguments of an existing function or member function to fixed values.\\
	\bigskip
	It appears extensively in \textbg{existing C++ and ROS 2 code}.
\end{frame}
\begin{frame}[fragile]{Lambdas and callbacks}{\texttt{std::function} and \texttt{std::bind}}
  \begin{columns}
		\column{.9\textwidth}
		\begin{lstlisting}[style=codebeamer, language=C++, caption=Storing heterogeneous callables in \texttt{std::function}.]
#include <functional>

int plain_fn(int x) { return x * 2; }

std::function<int(int)> f;
f = plain_fn;                    // a plain function
f = [](int x) { return x * 2; }; // a lambda

// std::function objects can be stored in containers
std::vector<std::function<void(float)>> clbks;
clbks.push_back([](float v) { std::cout << v << std::endl; });
clbks.push_back([](float v) { /* log v */ });
for (auto & cb : clbks) cb(3.14f);
    \end{lstlisting}
	\end{columns}
\end{frame}
\begin{frame}[fragile]{Lambdas and callbacks}{\texttt{std::function} and \texttt{std::bind}}
	\begin{columns}
		\column{.9\textwidth}
		\begin{lstlisting}[style=codebeamer, language=C++, caption=\texttt{std::bind} compared to its lambda equivalent.]
#include <functional>

int add(int a, int b) { return a + b; }

// Bind first argument to 10; second remains free (_1)
auto add10 = std::bind(&add, 10, std::placeholders::_1);
std::cout << add10(5) << std::endl; // 15

// The modern equivalent with a lambda
auto add10_lambda = [](int b) { return add(10, b); };
    \end{lstlisting}
	\end{columns}
\end{frame}
